\section{Solving the filled equations}
\label{sec:elliptic:construction}
\label{sec:elliptic}

The metric comparison and the energy estimate from
Section~\ref{sec:scalar:construction} become useful only after its two
scalar equations have a smooth solution.  We construct that solution
here, on the filled manifold and at fixed prepared data.  The boundary
of a large truncation is a single coordinate sphere; both unknowns
vanish there.

The first task is to keep the conformal variable above $-\epsilon$ and
to bound $Z$ without assuming uniform ellipticity.  The height comparison
and the floor calculation give the lower bound.  Testing the scalar
divergence equation on the graph of $f$ then gives a polynomial upper
bound for $Z$.  Together these bounds control the slope and make both
scalar equations uniformly elliptic for each fixed $N$.

We next prove the local estimates needed to pass from these bounds to
smoothness.  The gradient-aligned trace equation and the bounded
natural-growth divergence equation are treated separately, as reusable
scalar lemmas.  Finally we solve the trace equation for every trial
$Z$, use it to define a compact map, and obtain a fixed point by two
homotopies.  Exhausting the outer spheres gives the complete solution
and the end estimates needed for its ADM flux.  No regularity theorem
for elliptic systems or numerical Penrose inequality enters this
construction.

\begin{theorem}[Filled construction]\label{thm:elliptic:filled}
Fix smooth prepared data satisfying Equation~\eqref{eq:end:prepared}
on a complete one-ended exterior with nonempty compact smooth boundary.
In particular, $K$ is compactly supported,
$g-\delta=O_d(r^{2-n})$ for every fixed $d$,
$R_g=O(r^{-n-\beta})$ for some $0<\beta<1$,
the dominant-energy margin is bounded below by a positive multiple of
$\varrho^{-n-\beta}$, and every boundary component is strictly future
trapped.  Fix $0<\epsilon<1$ and
\[
 \frac n4<b_1<\min\{n-2,n/2\}.
\]
Use the filled manifolds $\mathcal M_N$, weights $\rho_0,\rho_1,\rho$,
and definitions in Equation~\eqref{eq:scalar:variables}.  Choose
$\delta_0>0$ and the polynomial $C_N$ as in
Lemma~\ref{lem:scalar:floor}.  For all sufficiently large integers
$N>k=n-1$, there are smooth functions $f,Z$ on $\mathcal M_N$
solving Equation~\eqref{eq:scalar:system}.  The associated function $t$
satisfies
\begin{equation}\label{eq:elliptic:polynomial-bounds}
 -\epsilon<t\le Z\le P(N),\qquad
 \ell\le l(t)\le e,\qquad
 |f|+|df|_g\le \exp(P(N)),\qquad \ell=e^{-N\epsilon}.
\end{equation}
Here and below $P$ is a polynomial with coefficients depending only on
the fixed prepared data, $n$, $\epsilon$, and the fixed choices in the
filling.  The height $h=\eta f$, $\eta=\ell^{3/2}$, obeys
\begin{equation}\label{eq:elliptic:height-bound}
 |h|\le C,\qquad |h(x)|\le C r(x)^{-b_1}
 \quad\hbox{on the original end},
\end{equation}
where $C$ is independent of $N$.  At each fixed $N$, for every integer
$j\ge0$ there are constants $C_{N,j}$ and $a_N>0$ such that on that end
\begin{equation}\label{eq:elliptic:end-bounds}
 |D^j f|\le C_{N,j}e^{-a_Nr},\qquad
 t,Z=O_j(r^{2-n}),\qquad
 \Delta_g t=O(r^{-n-\beta}).
\end{equation}
The constants in Equation~\eqref{eq:elliptic:end-bounds}, and higher
derivative bounds on compact sets, may depend arbitrarily on fixed $N$.
\end{theorem}

For clarity, we specify the homotopies used in the proof.  On a large
outer truncation $\mathcal M_{N,R}$ impose $f=Z=0$ on its sole boundary
$S_R$.  In the first homotopy replace $K$ everywhere by $K_a=aK$,
$0\le a\le1$, retaining the definitions of $\mathcal T,F,V$ with that
replacement, and put
\begin{equation}\label{eq:elliptic:homotopies}
 F=\eta f,\qquad
 \divg V=u\Xi_a,\qquad
 \Xi_a=\delta_0(\mathcal T+\eta a_w|df|)
       +\rho-m_0C_N(\rho_0+v\rho_1).
\end{equation}
The second homotopy keeps $K=0$ and replaces $\Xi_0$ by $s\Xi_0$,
$0\le s\le1$.  The weights and $C_N$ do not vary in either homotopy.
The endpoint $(a,s)=(0,1)$ is common to them.  Until the arbitrary-trial problem is introduced, the estimates concern
smooth fixed points on a truncation with $t\ge-\epsilon$.
The floor will be excluded on the boundary of the degree domain.
Solvability for arbitrary trials will be established separately, without
that restriction.

\subsection{Height comparison and the outer floor}

The trace equation controls the normalized height before any derivative
estimate.  Its end comparison also bounds the outer normal derivative;
this will prevent the floor from being reached on the truncating sphere.

\begin{lemma}\label{lem:elliptic:height}
There are constants $C$ and $r_1$, independent of $N,R$ and the homotopy
parameters, such that every smooth solution of the trace equation in
Equation~\eqref{eq:elliptic:homotopies} satisfies
\[
 |\eta f|\le C,\qquad
 |\eta f|\le C(r^{-b_1}-R^{-b_1})\quad(r_1\le r\le R).
\]
For each fixed $N$ there is $R_{\min}(N)$ such that $R\ge R_{\min}(N)$
and $Z=0$ on $S_R$ imply $t>-\epsilon$ there.  The same conclusions
hold in the second homotopy, where in fact $f=0$.
\end{lemma}

\begin{proof}
At a positive maximum of $f$, $df=0$, $A_\chi=I$, and
\[
 \eta f=\tr K_a+l\Delta_g f\le \tr K_a.
\]
The negative minimum gives the reverse bound.  The tensor fields in
the filling have a common pointwise bound, including along the product
necks, so this proves the first assertion uniformly.  If $K=0$, both
comparisons give $f=0$.

On the end $K_a=0$.  At a point of contact with
$H(r)=C(r^{-b_1}-R^{-b_1})$, considered as a test for $h=\eta f$, the
two gradients agree.  Thus the axial direction of $A_\chi$ is the
radial direction.  In the Euclidean metric the Hessian contraction is
\[
 A_\chi:\Hess H
 =Cb_1r^{-b_1-2}\big((b_1+1)\chi-k\big).
\]
The error from $g-\delta=O_2(r^{2-n})$ has absolute value at most
$C' r^{2-n}Cb_1r^{-b_1-2}$, uniformly for $0<\chi\le1$.
Since $b_1+1<k$, the contraction is strictly negative for $r\ge r_1$,
with $r_1$ independent of $N$.  The height equation is
\[
 A_\chi:\Hess h
   =\frac{\eta\sqrt D}{l}\,h.
\]
Consequently $H$ and $-H$ are respectively an upper and a lower
comparison function.  For $R\ge2r_1$, their constant $C$ can be chosen
to dominate the already established height bound on $S_{r_1}$.
The maximum principle then proves the displayed end estimate.

Taking the one-sided derivative at the outer zero boundary gives
\begin{equation}\label{eq:elliptic:outer-gradient}
 |df|_{S_R}\le C\eta^{-1}R^{-b_1-1}.
\end{equation}
Tangential derivatives vanish there.  For all real $t$ we have $l(t)\le e$.
The identity $Z=t+(2k)^{-1}\log(1+l(t)^2|df|^2)$ therefore gives, on
$S_R$,
\[
 0\ge t\ge-\frac1{2k}
       \log\bigl(1+e^2C^2\eta^{-2}R^{-2b_1-2}\bigr).
\]
Choosing $R_{\min}(N)$ sufficiently large proves the assertion.
\end{proof}

Together, Lemmas~\ref{lem:elliptic:height} and
\ref{lem:scalar:floor} exclude every fixed point with
$\min t=-\epsilon$.  The height lemma supplies the bounds assumed by
the scalar floor lemma and excludes the outer sphere by
Equation~\eqref{eq:elliptic:outer-gradient}.  The scalar floor lemma
then excludes an interior minimum in either homotopy.  These
exclusions concern the boundary of the degree domain; they make no
assertion about solutions whose minima lie strictly below $-\epsilon$.

\subsection{A polynomial upper bound before ellipticity}

Above the floor, the implicit relation expresses the metric distortion
directly in terms of $Z$.  We therefore need an upper bound for this
single scalar quantity.  The proof first bounds its high-level mass,
then combines a graph Sobolev inequality with exponential testing.
The bound for $Z$ must be polynomial in $N$: the enclosure argument
in Section~\ref{sec:height} needs metric distortion at most
$\exp(P(N))$ while $N$ increases.

\begin{proposition}\label{prop:elliptic:upper}
Every smooth fixed point on $\mathcal M_{N,R}$ with $t\ge-\epsilon$
satisfies Equation~\eqref{eq:elliptic:polynomial-bounds}, with the same
polynomial bounds for all $R\ge R_{\min}(N)$ and all homotopy parameters.
\end{proposition}

\begin{proof}
Write $\sigma=|df|$, $B=K_a(w,\cdot)$, $A=A_\chi$, and
\[
 d\mu_D=\sqrt D\,dV_g,\qquad
 W=e^{(k-1)t}l,\qquad u=W\sqrt D=l e^{kZ-t}.
\]
The floor already implies
\begin{equation}\label{eq:elliptic:elementary-distortion}
 \ell\le l\le e,\quad t\le Z,\quad
 D=e^{2k(Z-t)},\quad
 \sigma\le\ell^{-1}e^{k(Z+\epsilon)}.
\end{equation}
Only the upper bound for $Z$ remains to be proved.

\medskip\noindent\emph{High levels.}
On the support of $m_0$, the floor implies
$-\epsilon\le t\le-\epsilon+N^{-1}$ and $l\asymp\ell$.
In particular,
\[
 \eta a_w\sigma=\frac\eta l(\sqrt D-D^{-1/2}).
\]
The weights are bounded above uniformly, and $C_N$ is polynomial.
It follows from $D=e^{2k(Z-t)}$ that there is a positive polynomial
threshold $Z_0=P(N)$ such that
\begin{equation}\label{eq:elliptic:high-source}
 \Xi_a\ge\delta_0\mathcal T+\rho
                   +\frac{\delta_0}2\eta a_w\sigma
 \quad\hbox{when }Z>Z_0.
\end{equation}
For example, it suffices that
$e^{kZ_0}\ge C(1+C_N)\eta^{-1}$, after increasing the fixed constant;
its logarithm has polynomial size.  Away from the support of $m_0$
the inequality is immediate.

Choose a smooth nondecreasing $\zeta$ that vanishes on $(-\infty,Z_0]$
and equals one on $[Z_0+1,\infty)$, with $0\le\zeta'\le C$.
Testing the divergence equation with $\zeta(Z)$ gives
\[
 \int \zeta(Z)u\Xi_a
  =-\int\zeta'(Z)u\{k|dZ|_A^2+\langle B,dZ\rangle_A\}
  \le C\int_{\{Z_0<Z<Z_0+1\}\cap\supp K}u|K|^2.
\]
The integration has no boundary term since $Z=0<Z_0$ on $S_R$.
On the last strip $u=l e^{kZ-t}\le\exp(P(N))$.
The support of $K$ consists of a fixed exterior compact set and the
enlarged filling; it has volume $O(1+N)$.  Thus
\begin{equation}\label{eq:elliptic:high-integral}
 \int_{\{Z>Z_0+1\}}u
       (\delta_0\mathcal T+\rho+\tfrac12\delta_0\eta a_w\sigma)
       \,dV_g\le\exp(P(N)).
\end{equation}
This estimate uses no upper bound for $\sigma$.  It controls the high
levels of $Z$ with exactly the volume and slope weights needed on the
graph.  We next convert it into the initial integrability for an
iteration on fixed base patches.

Fix $0<\alpha\le k-1$, independent of $N$.  Cover the enlarged core
and a fixed extra end annulus by base coordinate patches of uniformly
controlled size and geometry.  The number of patches is $O(1+N)$,
and $\rho$ has a uniform positive lower bound on them.  We claim that
on each such patch $Q$,
\begin{equation}\label{eq:elliptic:initial-integral}
 \int_Q e^{\alpha Z}\,d\mu_D\le\exp(P(N)).
\end{equation}
Below $Z_0+1$, Equation~\eqref{eq:elliptic:elementary-distortion}
proves this directly.  Above that level, when $D\le2$ use
\[
 \frac{e^{\alpha Z}\sqrt D}{u}
   =l^{-1}e^{(\alpha-k+1)t}D^{\alpha/(2k)}
   \le\exp(P(N)).
\]
When $D\ge2$, use instead
\[
 u\eta a_w\sigma=\eta e^{(k-1)t}(D-1),\qquad
 e^{\alpha Z}\sqrt D=e^{\alpha t}D^{1/2+\alpha/(2k)}.
\]
The ratio of the second expression to the first is at most
$C\eta^{-1}e^{(\alpha-k+1)t}\le\exp(P(N))$ because
$1/2+\alpha/(2k)\le1$ and $t\ge-\epsilon$.
Equation~\eqref{eq:elliptic:high-integral} proves the claim.

\medskip\noindent\emph{Sobolev inequality on the product graph.}
We use the ordinary graph of $f$ in $(Q\times\R,g+dz^2)$ so that
its ambient geometry depends only on the fixed background patch.  The
warp $l(t)$ has no a priori derivative bound yet.  We compare the
ordinary graph's measure and inverse metric with $d\mu_D$ and $A$,
and use $\mathcal T$ to control its mean curvature.
The ordinary graph has measure and inverse metric
\[
 \sqrt{1+\sigma^2}\,dV_g,
 \qquad A_* = I-\frac{df\otimes df}{1+\sigma^2}.
\]
The ratios of $D$ and $1+\sigma^2$ lie between $\exp(-P(N))$ and
$\exp(P(N))$, even when $\sigma$ is unbounded.  The axial eigenvalues
$d$ and $\chi=kd/(k+pv)$ differ by at most a factor $1+N/k$.
Thus these graph norms compare with $d\mu_D$ and $A$ with factors
$\exp(P(N))$.

Its mean curvature is
$(1+\sigma^2)^{-1/2}\tr_{A_*}\Hess f$.  Write
$K_a+H^f$ in its transverse and axial blocks.  The transverse
contraction is controlled by $|T|$, and the axial coefficient of
$A_*$ is at most $e^2d$.  Also
$\sqrt D/(l\sqrt{1+\sigma^2})\le\exp(P(N))$.
Proposition~\ref{prop:scalar:coercivity} therefore gives
\[
 |H_{\mathrm{graph}}|\le\exp(P(N))(1+\sqrt{\mathcal T}).
\]
Extend the metric of a slightly larger fixed base patch to a closed
manifold and use Nash's smooth isometric embedding theorem
\cite{Nash1956}; then take its product with the line.  The additional second
fundamental form is uniformly bounded.  Such embeddings can be chosen
from finitely many fixed models for the cap, transitions, and product
necks. The Michael--Simon inequality~\cite{MichaelSimon1973}, in the
formulation of \cite[Chapter~4, Section~6, Theorem~6.7]{SimonGMT},
applied to compactly supported functions on this smooth $n$-dimensional graph,
then gives, after its usual $L^1$ to $L^2$ conversion,
\begin{equation}\label{eq:elliptic:graph-sobolev}
 \left(\int |\varphi|^{2\kappa_n}\,d\mu_D\right)^{1/\kappa_n}
 \le\exp(P(N))\int
       \bigl(|d\varphi|_A^2+(1+\mathcal T)\varphi^2\bigr)\,d\mu_D,
 \qquad \kappa_n=\frac n{n-2}.
\end{equation}
To see the conversion explicitly, apply the $L^1$ inequality to
$|\varphi|^{2(n-1)/(n-2)}$ and use Cauchy--Schwarz on its derivative
and mean-curvature terms.  The resulting common
$L^{2n/(n-2)}$ factor cancels.  No bound for the height or area of the
graph is required.

\medskip\noindent\emph{Exponential testing.}
The graph inequality still contains $\mathcal T$ on its right side.
The divergence equation controls that term together with the gradient
of an exponential of $Z$; this is the estimate that closes the iteration.
Globally $\Xi_a\ge\delta_0\mathcal T-P(N)$.
Let $\psi$ be a base cutoff compactly supported in $Q$ and test the
equation with $\psi^2e^{2qZ}/W$.  Put
$\beta_W=d\log W=(k-1+p)dt$.
Since $A\le A_d$, scalar coercivity gives
$|\beta_W|_A^2\le P(N)\mathcal T$.
Integration by parts yields exactly
\begin{align*}
 \int \psi^2e^{2qZ}\Xi_a\,d\mu_D
 =-\int e^{2qZ}
  \bigl(2q\psi^2dZ-\psi^2\beta_W+2\psi d\psi\bigr)
                   \cdot A(kdZ+B)\,d\mu_D.
\end{align*}
Young's inequality bounds the $\beta_W\cdot A\,dZ$ term by
$\delta_0\mathcal T/8+P(N)|dZ|_A^2$, the $\beta_W\cdot AB$
term by $\delta_0\mathcal T/8+P(N)$, and the $qB\cdot A\,dZ$
term by a fixed fraction of $kq|dZ|_A^2+C q$.
The cutoff terms cost a further fixed fraction of that gradient term
and $P(N)(\psi^2+|d\psi|^2)$.
Consequently there is a polynomial $q_*(N)\ge\max\{1,\alpha\}$ such
that, for $q\ge q_*(N)$,
\begin{equation}\label{eq:elliptic:exponential-energy}
 \int e^{2qZ}\psi^2
           (\mathcal T+q|dZ|_A^2)\,d\mu_D
 \le P(N)(1+q)\int e^{2qZ}
                    (\psi^2+|d\psi|_g^2)\,d\mu_D.
\end{equation}
Only $dt$ occurs in this test; it does not differentiate $p$.

Apply Equation~\eqref{eq:elliptic:graph-sobolev} with
$\varphi=\psi e^{qZ}$ and then use
Equation~\eqref{eq:elliptic:exponential-energy}.  For nested base
patches $Q'\Subset Q$ separated by a coordinate distance $b\le1$,
this gives
\[
 \|e^Z\|_{2\kappa_nq,Q'}
 \le\bigl[e^{P(N)}(1+q)^2b^{-2}\bigr]^{1/(2q)}
                          \|e^Z\|_{2q,Q}.
\]
All norms here use $d\mu_D$.  Iterate with
$q_i=\kappa_n^iq_*$ and geometrically decreasing patch gaps.
The sums of $q_i^{-1}$ and $i q_i^{-1}$ converge, so
\begin{equation}\label{eq:elliptic:moser-output}
 \sup_{Q'}e^Z
 \le e^{P(N)}b^{-C}\|e^Z\|_{2q_*,Q},
\end{equation}
where $C$ is independent of $N$.

Finally interpolate the norm on the right with
Equation~\eqref{eq:elliptic:initial-integral}:
\[
 \|e^Z\|_{2q_*,Q}
 \le \bigl(\sup_Qe^Z\bigr)^\theta e^{P(N)},
 \qquad \theta=1-\frac\alpha{2q_*}<1.
\]
Choose expanding patches $Q_j$, all inside one fixed larger patch,
with gaps comparable to $2^{-j}$.  The logarithms $M_j=\log\sup_{Q_j}e^Z$
satisfy $M_j\le\theta M_{j+1}+P(N)+Cj$.
Iteration gives
\[
 M_0\le\theta^mM_m+
          \sum_{j=0}^{m-1}\theta^j(P(N)+Cj)\le\theta^mM_m+P(N).
\]
The last polynomial is justified by
$(1-\theta)^{-1}=2q_*/\alpha\le P(N)$.
For each individual smooth solution the suprema on the largest patch
are finite, so the first term tends to zero.  We have proved
$Z\le P(N)$ on the enlarged core and the drift support.  A larger
maximum outside that set would have $K_a=0$, $dZ=0$, and, by
Equation~\eqref{eq:elliptic:high-source}, $\Xi_a>0$; at the same point
$u^{-1}\divg(kuA\nabla Z)=kA:\Hess Z\le0$, a contradiction.

For the second homotopy $f=0$.  A positive maximum above the penalty
band has right side $s\rho>0$ if $s>0$, which is impossible.  When
$s=0$ the solution is zero.  This proves its upper bound as well.
Equation~\eqref{eq:elliptic:elementary-distortion} and the height bound
prove the remaining bounds in
Equation~\eqref{eq:elliptic:polynomial-bounds}.  In particular,
\[
 d\ge e^{-2k(P(N)+\epsilon)},\qquad
 \chi\ge\frac{k}{k+N}e^{-2k(P(N)+\epsilon)}.
\]
The same bounds give positive upper and lower bounds for $u$ at each
fixed $N$.  Thus both $A_\chi$ and $kuA_\chi$ are uniformly elliptic,
with constants independent of the truncating radius and the homotopy
parameter.  These are the ellipticity bounds used in the following
scalar estimates.
\end{proof}

\subsection{A scalar gradient estimate with a measurable axial coefficient}

The upper bound has made both scalar equations uniformly elliptic
at each fixed $N$.  Regularity still has to pass between them in a
specific order.  Dividing the trace equation by $l/\sqrt D$ gives
\[
 A_\chi:\Hess_g f
   =\frac{\sqrt D}{l}\bigl(\eta f-\tr_{A_\chi}K_a\bigr).
\]
Its right side and $Df$ are bounded, but its coefficients depend on
$Z$, for which we have no continuity modulus.  Schauder theory does
not yet apply.

The gradient direction accounts for the matrix structure; the remaining
unknown eigenvalue $\chi$ can be treated as merely measurable.  The next
lemma gives both a H\"older bound for $Df$ and a local mass bound for
$|D^2f|^2$.  This second conclusion is needed because the divergence
source for $Z$ contains $|D^2f|^2$: the mass bound will give H\"older
continuity of $Z$ through the following scalar lemma.  Once both $Z$
and $Df$ have a modulus, scalar Schauder estimates can begin.  All
constants in this regularity argument may depend arbitrarily on fixed
$N$; the bounds needed as $N$ increases were established above.

We prove the gradient estimate in two regimes.  Strict rank-one
absorption first gives $W^{2,s}$ control for some $s>2$.  A
trace-reversed Hessian flux then shows that either the gradient norm
drops on a smaller ball or the solution is close to an affine function
with nonzero gradient.  In the latter regime we freeze the direction
while retaining the measurable axial coefficient; the axial derivative
satisfies a scalar divergence equation.  Iterating these two alternatives
gives gradient continuity and the required Hessian mass bound.

\begin{lemma}[Gradient-aligned rank-one equation]
\label{lem:elliptic:gradient}
Let $n\geq3$ and $0<c\leq1$.  Let $Q$ be a smooth metric coordinate ball, or a
smooth half-ball with constant Dirichlet data on its face.  Suppose that
$u$ is $C^2$ up to the face and satisfies almost everywhere
\begin{equation}
 A:\Hess_g u=G,\qquad
 A=\id-(1-\chi)e\otimes e,\qquad
 c\leq\chi\leq1,\qquad
 e=\frac{\nabla_g u}{|\nabla_g u|_g}
 \quad\text{where }\nabla_g u\ne0.
 \label{eq:elliptic:rankone-equation}
\end{equation}
At a zero of the gradient, $e$ may be any measurable unit vector.  Assume
$|\nabla_g u|_g+|G|\leq M$.  On every smaller ball, including its portion
of the Dirichlet face, there are estimates
\begin{equation}
 [Du]_{C^\alpha}\leq C,\qquad
 \int_{B_r(x)\cap Q}|\Hess_g u|_g^2\,dV_g
       \leq C r^{n-2+2\alpha}.
 \label{eq:elliptic:rankone-conclusion}
\end{equation}
Here $\alpha\in(0,1)$ depends only on $n,c$, and $C$ depends on $n,c,M$,
the fixed coordinate geometry and the distance from the other chart
boundaries.  No modulus of continuity of $\chi$ enters these constants.
\end{lemma}

\begin{proof}
We use only the scalar linear De Giorgi--Nash and weak Harnack estimates,
the Krylov--Safonov interior estimate for strong nondivergence solutions,
and the Euclidean Hessian estimates of Calder\'on and Zygmund;
see~\cite{GilbargTrudinger}.  The two structural
arguments needed to apply these results are proved below.  All coordinate
norms of matrices use the Frobenius norm unless indicated otherwise.

\medskip\noindent
\emph{1. A Hessian estimate with a strict absorption margin.}
Normalize the metric at the center of a chart.  The coordinate principal
matrix is
\[
 a^{ij}=g^{ij}-(1-\chi)e^i e^j.
\]
If $g=I$, then $|I-a|=1-\chi\leq1-c$.  Consequently, on a sufficiently
small chart, depending only on $c$ and the metric modulus,
\begin{equation}
 |I-a|\leq1-\frac c2.
 \label{eq:elliptic:cordes-margin}
\end{equation}
This remains true for the reflected metrics used below, since they are
continuous and uniformly Lipschitz.  In particular this assertion places
no condition on the oscillation of $e$ or $\chi$.

Let $\mathcal R=D^2\Delta^{-1}$ on $\R^n$, with its matrix-valued range
given the Frobenius norm.  Plancherel's identity gives
$\|\mathcal R\|_{L^2\to L^2}=1$.  Its $L^4$ norm is finite by the
Calder\'on--Zygmund theorem.  Interpolation therefore supplies
$s_0>2$, depending only on $n,c$, such that
\[
 \|\mathcal R\|_{L^s\to L^s}(1-c/2)<1,
 \qquad 2\leq s\leq s_0.
\]
For a compactly supported $v$ the identity
$\Delta v=a:D^2v+(I-a):D^2v$ now gives, by absorption,
\[
 \|D^2v\|_{L^s}\leq C\|a:D^2v\|_{L^s},
 \qquad s=2,s_0.
\]
Cutoffs, first-derivative interpolation and absorption on nested balls
give the local version
\begin{equation}
 \|D^2v\|_{L^s(B_{r/2})}
 \leq C\bigl(\|a:D^2v\|_{L^s(B_r)}
              +r^{-2}\|v\|_{L^s(B_r)}\bigr),
 \qquad s=2,s_0.
 \label{eq:elliptic:cordes-local}
\end{equation}
For completeness, let $\psi$ be a smooth cutoff for the nested balls.
Applying the compact-support estimate to $\psi v$ first produces an additional
$(t-\rho)^{-1}\|Dv\|_{L^s(B_t)}$ on $B_\rho\Subset B_t$.
The interpolation estimate on a slightly larger ball bounds this by
$\varepsilon\|D^2v\|_{L^s}$ plus
$C\varepsilon^{-1}(t-\rho)^{-2}\|v\|_{L^s}$, with the scale factors
absorbed into $\varepsilon$.  Choose the coefficient of the first term
smaller than the square of the geometric ratio of the successive gaps,
and sum over increasing nested radii.  The resulting geometric series
is exactly Equation~\eqref{eq:elliptic:cordes-local}.
These estimates initially apply to the individually finite Sobolev norms
of the functions under consideration and then, by approximation, to
$W^{2,s}$ functions.

\medskip\noindent
\emph{2. The exact flux and the constant-value reflection.}
Write $P=\nabla_g u$, $H=\Hess_g u$, and $q=|P|_g^2/2$.  The equation
implies
$\Delta_g u=(1-\chi)H(e,e)+G$.  Since
$\nabla_g q=H^\sharp P$, it gives the pointwise identity
\begin{equation}
 A\nabla_g q-GP
      =\bigl(H^\sharp-(\Delta_g u)\id\bigr)P=:\mathcal F.
 \label{eq:elliptic:trace-reversed-flux}
\end{equation}
Both sides vanish when $P=0$, so no choice of axis at those points affects
this identity.  For smooth $u$, commuting third derivatives gives
\begin{equation}
 \divg_g\mathcal F
     =|H|_g^2-(\Delta_g u)^2+\Ric_g(P,P).
\label{eq:elliptic:trace-reversed-divergence}
\end{equation}
For $C^2$ functions the same formula holds distributionally: approximate
$u$ in $C^2$ on compact subsets of each original smooth side.  The flux
and the right side converge uniformly, so integration against a smooth
compactly supported test function passes to the limit.  This argument
establishes only the geometric divergence identity; the algebraic
Equation~\eqref{eq:elliptic:trace-reversed-flux} is applied to the original
solution, and no approximant is required to solve the equation.
Neither identity differentiates $\chi$.

Put $\gamma=1-(1-c)^2>0$.  Young's inequality, with the gap $\gamma/2$,
gives
\[
 ((1-\chi)H(e,e)+G)^2
     \leq(1-\gamma/2)|H|_g^2+C_cG^2.
\]
Thus, when $|P|_g\leq1$ and $s=1-|P|_g^2\geq0$,
\begin{equation}
 \divg_g(A\nabla_g s+2GP)
       \leq-\gamma|H|_g^2+C_cG^2+2|\Ric_g|_g.
 \label{eq:elliptic:gradient-deficit}
\end{equation}

We record precisely how this inequality is used at the face.  Subtract
the constant boundary value and take Gaussian coordinates
$g=dt^2+h_{ab}(y,t)\,dy^a dy^b$ on $t\geq0$.  Extend $u$ oddly and
$h$ evenly across $t=0$, and extend $\chi$ evenly and $G$ oddly.
The reflected equation holds almost everywhere.  Tangential derivatives
of $u$ vanish on the face, while the normal derivative agrees from both
sides.  Therefore $u$ is $C^1$ after reflection and belongs locally to
$W^{2,s}$ for every finite $s$ individually.  The metric is Lipschitz,
smooth on either side, and $s$ is a continuous $W^{1,2}$ function.

Let $J=\sqrt{\det g}$ and define the face mean curvature using the normal
$+\partial_t$ by
\[
 \kappa(y)=\frac12 h^{ab}(y,0)\partial_t h_{ab}(y,0+),
 \qquad p(y)=u_t(y,0).
\]
On the upper side, $H_{ab}=\tfrac12(\partial_t h_{ab})p$ and hence
$\mathcal F^t_+=-\kappa p^2$.  On the lower side the signs of these
tangential Hessians reverse and $\mathcal F^t_- =\kappa p^2$.
Since $J$ is continuous across the face, the density-valued flux
$J(A\nabla_gs+2GP)=-2J\mathcal F$ has normal jump
\[
 j(y)=4J(y,0)\kappa(y)p(y)^2.
\]
Adding the coordinate vector
\begin{equation}
 B_{\rm face}(y,t)=-j(y)\mathbf1_{\{t>0\}}\partial_t
 \label{eq:elliptic:face-step}
\end{equation}
cancels that interface measure exactly.  Indeed its coordinate
divergence is $-j(y)\delta_{\{t=0\}}$; no derivative in $y$ occurs,
because the vector has only a $t$ component.  Its $L^\infty$ norm is
at most $C\|\kappa\|_\infty$ under the gradient normalization.
All these statements are local, so the coefficient $j(y)$ only needs
to be extended constantly in the normal direction over the chart.

On a normalized unit ball let the metric oscillation, Christoffel symbols,
smooth-side curvature, face second fundamental form if present, and
forcing be at most $\varepsilon$.  The reflected and interior cases
then both give the distributional inequality
\begin{equation}
 \divg(a_1Ds+B_1)
       \leq-\gamma_1|H|_g^2+C\varepsilon,
 \qquad
 a_1^{ij}=Ja^{ij},\qquad
 \|B_1\|_\infty\leq C\varepsilon,
 \label{eq:elliptic:deficit-coordinate}
\end{equation}
where $B_1^i=2JG P^i+B_{\rm face}^i$ and $\gamma_1>0$ is fixed.
The matrix $a_1$ is symmetric and uniformly elliptic with fixed constants.
The face can be absent or can be a plane at any position in the ball.

\medskip\noindent
\emph{3. Gradient-norm drop or approximation by a nonzero affine function.}
Fix $r_*\in(0,1/32)$.  For every $b_*>0$ there are
$k_*\in(r_*,1)$ and $\varepsilon_*>0$ with the following property.
Suppose that $|\nabla_g u|_g\leq1$ on $B_1$, the normalized geometric
and forcing errors just listed are at most $\varepsilon_*$, and
$u(0)=0$.  Then either
\begin{equation}
 \sup_{B_{r_*}}|\nabla_g u|_g\leq k_*,
 \label{eq:elliptic:norm-drop}
\end{equation}
or there is an affine $L$ such that
\begin{equation}
 \frac12\leq|DL|\leq2,\qquad
 \sup_{B_{r_*}}|u-L|\leq b_*r_*.
 \label{eq:elliptic:affine-entry}
\end{equation}

Here is the compactness argument with its energy step.  Otherwise choose
a sequence with errors tending to zero, no approximation in
Equation~\eqref{eq:elliptic:affine-entry}, and gradient supremum tending
to one on $B_{r_*}$.  Its functions are uniformly Lipschitz and vanish
at zero.  Equation~\eqref{eq:elliptic:cordes-local}, applied to
\[
 a:D^2u=G+a^{ij}\Gamma^\ell_{ij}u_\ell,
\]
gives uniform local $L^2$ bounds for $D^2u$, hence for $H$ and $Ds$.
The weak Harnack inequality applies to the nonnegative weak
supersolution $s$ in Equation~\eqref{eq:elliptic:deficit-coordinate}
after discarding the negative Hessian term.  It gives, for some $p_0>0$,
\[
 \left(\frac1{|B_{1/3}|}\int_{B_{1/3}}s^{p_0}\right)^{1/p_0}
   \leq C\left(\inf_{B_{1/3}}s+\|B_1\|_\infty+\varepsilon\right)
   \longrightarrow0.
\]
The infimum tends to zero because $B_{r_*}\subset B_{1/3}$ and the
gradient supremum there tends to one.  Thus $s\to0$ in measure locally.

To recover the Hessian energy, take a compactly supported cutoff $\psi$
in $B_{1/3}$ and test the same inequality with
$\psi^2(b-s)_+$, where $0<b<1$ is fixed.  Its weak form is
\[
 \int(a_1Ds+B_1)\cdot D\varphi
    \geq\gamma_1\int|H|_g^2\varphi-C\varepsilon\int\varphi,
 \qquad\varphi\geq0.
\]
Dropping the nonnegative Hessian term, expanding the test derivative and
using Young's inequality proves
\begin{equation}
 \int_{\{s<b\}}\psi^2|Ds|^2
   \leq C_\psi b^2+C_\psi\|B_1\|_\infty^2
          +C_\psi b\|B_1\|_\infty+C_\psi\varepsilon b.
 \label{eq:elliptic:lower-truncation}
\end{equation}
On a fixed compact subset, the integral of $|Ds|$ over $\{s<b\}$ has
limsup at most $Cb$.  Its integral over $\{s\geq b\}$ tends to zero
by the uniform $L^2$ bound for $Ds$ and convergence in measure of $s$.
Letting $b\downarrow0$ yields $Ds\to0$ locally in $L^1$.
Testing Equation~\eqref{eq:elliptic:deficit-coordinate} with a fixed
nonnegative smooth cutoff now yields $H\to0$ locally in $L^2$.
Since $\Gamma\to0$ and $Du$ is bounded, also $D^2u\to0$ locally in
$L^2$.

Uniform Lipschitz compactness produces a locally uniform limit.  Its
distributional Hessian vanishes, so it is affine.  Poincar\'e's inequality
and $D^2u\to0$ give convergence of the gradients to its slope in local
$L^2$.  Since $s\to0$ in measure and $g\to I$, that slope has norm one.
This contradicts the failure of Equation~\eqref{eq:elliptic:affine-entry}
and proves the dichotomy.  Increasing $k_*$ if necessary ensures
$r_*<k_*<1$.

Repeated gradient-norm drops will give decay directly.  To handle the
other alternative, we next improve approximation by a nonzero affine
function.  Its fixed slope determines the direction to freeze; the
axial eigenvalue remains measurable throughout this argument.

\medskip\noindent
\emph{4. The frozen-axis linear equation.}
Let $e_0$ be a fixed Euclidean unit vector.  We first prove the quantitative
statement
\begin{equation}
 \|Y_0\|_{C^{1,\alpha_1}(B_{1/2})}
      \leq C\|Y_0\|_{W^{2,2}(B_1)}
 \quad\text{if}\quad
 \Delta_{e_0^\perp}Y_0+\chi\partial_{e_0}^2Y_0=0
 \ \text{a.e. in }B_1,
 \label{eq:elliptic:frozen-axis-estimate}
\end{equation}
for $Y_0\in W^{2,2}$ and arbitrary measurable $c\leq\chi\leq1$.
Rotate so that $e_0=\partial_n$ and put $w=\partial_nY_0\in W^{1,2}$.
Differentiating the displayed equation only in the sense of distributions
gives exactly
\begin{equation}
 \sum_{a<n}\partial_a^2w+\partial_n(\chi\partial_nw)=0,
 \quad\text{or}\quad
 \divg\bigl((I-e_0\otimes e_0+\chi e_0\otimes e_0)Dw\bigr)=0.
 \label{eq:elliptic:frozen-derivative-divergence}
\end{equation}
This is a scalar divergence equation for a $W^{1,2}$ weak solution with
a measurable uniformly elliptic matrix.  The De Giorgi--Nash estimate
gives $w\in C^{\alpha_1}$ on smaller balls, for a fixed
$0<\alpha_1<1$, quantitatively in its $L^2$ norm.  Testing with
$\psi^2(w-w(x))$, where $\psi$ is a cutoff on the ball, and using that
H\"older bound gives, uniformly for balls in a smaller fixed ball,
\begin{equation}
 \int_{B_r(x)}|Dw|^2\leq C r^{n-2+2\alpha_1}.
 \label{eq:elliptic:frozen-derivative-morrey}
\end{equation}
The ordinary Laplacian of $Y_0$ is
$f=(1-\chi)\partial_nw$.  By Cauchy--Schwarz,
\begin{equation}
 \int_{B_r(x)}|f|\leq C r^{n-1+\alpha_1}.
 \label{eq:elliptic:frozen-source-mass}
\end{equation}

Here is why this controls every component of $DY_0$.  Cut $f$ off on
an intermediate ball and form its Newton potential $V$.  For two points
at distance $r$ in a smaller ball, the part of $DV$ within distance
$2r$ of either point is bounded by
\[
 C\sum_{j\geq0}(2^{-j}r)^{1-n}(2^{-j}r)^{n-1+\alpha_1}
       \leq Cr^{\alpha_1}.
\]
For distances $R\geq2r$, the difference of the two gradient kernels is
at most $CrR^{-n}$.  Summing over dyadic annuli and using
Equation~\eqref{eq:elliptic:frozen-source-mass} gives
\[
 Cr\sum_{r\leq R\leq1}R^{\alpha_1-1}
       \leq Cr^{\alpha_1}.
\]
The farther part is bounded by the total $L^1$ norm of $f$ and contributes
$O(r)$.  The same estimates show that $V,DV$ are bounded.  The difference
$Y_0-V$ is harmonic on the smaller ball, so interior harmonic estimates
control its gradient and its gradient H\"older norm by its $L^2$ norm.
Together with the initial $W^{2,2}$ bound this proves
Equation~\eqref{eq:elliptic:frozen-axis-estimate}.  Taking slightly nested
balls in this argument yields the stated $B_{1/2}$ estimate.

\medskip\noindent
\emph{5. Improvement near a nonzero affine function.}
Fix a universal gradient bound, say $|\nabla_g u|_g\leq4$.
There are $\alpha_2\in(0,\alpha_1)$, $\lambda\in(0,1/8)$ and small
$b_0,\delta>0$, depending only on $n,c$, with the following property.
If $L$ is affine and
\begin{equation}
 \frac14\leq|DL|\leq4,\qquad
 \|u-L\|_{L^\infty(B_1)}\leq b,\quad 0<b\leq b_0,
 \qquad
 \|g-I\|_\infty+\|\Gamma\|_\infty+\|G\|_\infty\leq b\delta,
 \label{eq:elliptic:affine-improvement-hypotheses}
\end{equation}
then there is an affine $L'$ such that
\begin{equation}
 \sup_{B_\lambda}|u-L'|\leq b\lambda^{1+\alpha_2},
 \qquad |DL'-DL|\leq Cb.
 \label{eq:elliptic:affine-improvement}
\end{equation}
The metrics may again be smooth or reflected.

Set $Y=(u-L)/b$.  Its coordinate equation is
\[
 a:D^2Y=b^{-1}\bigl(G+a^{ij}\Gamma^\ell_{ij}u_\ell\bigr),
\]
whose right side has absolute value at most $C\delta$.
Equation~\eqref{eq:elliptic:cordes-local} and interpolation give
\begin{equation}
 \|Y\|_{W^{2,s_0}(B_{7/8})}\leq C.
 \label{eq:elliptic:normalized-remainder-sobolev}
\end{equation}
Put $e_0=DL/|DL|$ and $a_0=I-(1-\chi)e_0\otimes e_0$.
Since $Du=DL+bDY$ and $|DL|\geq1/4$, the direction projections obey
\[
 \|a-a_0\|_{L^{s_0}(B_{7/8})}\leq Cb,
 \qquad\|a-a_0\|_\infty\leq C.
\]
Indeed the map taking a nonzero vector to its unit-direction projection
has the estimate
$|\Pi(v)-\Pi(q)|\leq C\min(1,|v-q|/|q|)$ when $q\ne0$.
At $v=0$ any assigned unit projection satisfies the same estimate after
enlarging $C$.  The change from $Du$ to its metric gradient contributes
$O(\|g-I\|_\infty)$.
H\"older's inequality, with
$1/2=1/s_0+1/p$ and interpolation of the bounded matrix difference in
$L^p$, therefore gives
\begin{equation}
 R:=a_0:D^2Y,\qquad
 \|R\|_{L^2(B_{3/4})}\leq C(\delta+b^\sigma),
 \quad \sigma=\min\{1,(s_0-2)/2\}>0.
 \label{eq:elliptic:frozen-residual}
\end{equation}

We remove this residual in $W^{2,2}$, without any Sobolev embedding into
$C^0$.  On the convex ball $\mathcal B=B_{3/4}$, integration by parts
for a zero-trace function gives
\[
 \int_{\mathcal B}|D^2v|^2
   \leq\int_{\mathcal B}|\Delta v|^2.
\]
The omitted boundary term is the nonnegative integral of the sphere's
mean curvature times $(\partial_\nu v)^2$; the identity extends by
density to $W^{2,2}\cap W^{1,2}_0$.
Thus the Dirichlet solution operator $\Delta_D^{-1}$ satisfies
$\|D^2\Delta_D^{-1}f\|_2\leq\|f\|_2$.  The map
\[
 v\longmapsto\Delta_D^{-1}\bigl(R+(1-\chi)\partial_{e_0}^2v\bigr)
\]
is a contraction, with factor at most $1-c$, on
$W^{2,2}(\mathcal B)\cap W^{1,2}_0(\mathcal B)$ equipped with the Hessian
norm.  Its fixed point satisfies
\[
 a_0:D^2v=R,\qquad
 \|v\|_{W^{2,2}(\mathcal B)}\leq C\|R\|_2.
\]
Consequently $Y_0=Y-v$ satisfies the homogeneous frozen equation and has
bounded $W^{2,2}(\mathcal B)$ norm.  By
Equation~\eqref{eq:elliptic:frozen-axis-estimate},
$\|Y_0\|_{C^{1,\alpha_1}(B_{3/8})}\leq C$.

Independently, $Y$ has a uniform $C^\theta$ bound on $B_{1/2}$ by the
Krylov--Safonov estimate applied to its unfrozen equation: the matrix
$a$ is measurable and uniformly elliptic, its source is bounded, and
$\|Y\|_\infty\leq1$.  The function $Y$ is individually $W^{2,n}$
(also across the face), so the strong-solution version of that theorem
applies.  Hence $v=Y-Y_0$ has uniformly bounded $C^\theta$ norm on
$B_{3/8}$.  Its small $L^2$ norm now implies
\begin{equation}
 \|v\|_{L^\infty(B_{1/3})}
  \leq C\|v\|_2^{2\theta/(n+2\theta)}
  \leq C(\delta+b^\sigma)^{2\theta/(n+2\theta)}.
 \label{eq:elliptic:correction-uniform}
\end{equation}
To see the first inequality for small $\|v\|_2$, a value $|v(x)|=m$
forces $|v|\geq m/2$ on a ball of radius $(m/(2C))^{1/\theta}$;
integrating its square gives the stated exponent.

Choose $\alpha_2<\alpha_1$, then $\lambda$ so small that
$C\lambda^{1+\alpha_1}\leq\tfrac12\lambda^{1+\alpha_2}$, and finally
$\delta,b_0$ so small that the right side of
Equation~\eqref{eq:elliptic:correction-uniform} is at most
$\tfrac12\lambda^{1+\alpha_2}$.  The affine function
\[
 L'=L+b\bigl(Y_0(0)+DY_0(0)\cdot x\bigr)
\]
then satisfies Equation~\eqref{eq:elliptic:affine-improvement}.

\medskip\noindent
\emph{6. Iteration, including its transition between the two alternatives.}
Decrease $b_0$ further so that
$Cb_0/(1-\lambda^{\alpha_2})\leq1/8$, where $C$ is the slope constant
in Equation~\eqref{eq:elliptic:affine-improvement}.  Apply the dichotomy
with $b_*=b_0$.  At each center, first normalize $g$ there and choose a
fixed sufficiently small radius $R_0$.  Divide heights by $K R_0$, where
$K\geq\max(1,\|\nabla_g u\|_\infty)$, and subtract the center value.
On the rescaled unit ball the gradient is at most one.  The forcing is
multiplied by $R_0/K$, the connection and face second form by $R_0$,
and smooth-side curvature by $R_0^2$.  Metric oscillation from its center
value is $O(R_0)$.  Choose $R_0$ uniformly so that these errors satisfy
both the dichotomy threshold and $b_0\delta$.

Denote this normalized function by $u_0$.  After $j$ successive norm
drops, the normalized function is
\[
 u_j(x)=\frac{u_0(r_*^j x)}{r_*^j k_*^j}.
\]
It again has gradient at most one.  Its forcing has acquired the factor
$(r_*/k_*)^j$, and its metric, connection and face errors have acquired
at least the factor $r_*^j$.  All required smallness conditions persist
because $r_*<k_*$.  If this branch continues forever, then
\[
 \sup_{B_{r_*^j}}|u_0-u_0(0)|
       \leq C r_*^j k_*^j
       =C(r_*^j)^{1+\alpha_d},
 \qquad\alpha_d=\frac{\log k_*}{\log r_*}\in(0,1).
\]

If the affine alternative first occurs after $j$ norm drops, rescale its
ball once by $r_*$, retaining its current gradient normalization.
The resulting function $U$ has a unit-ball affine error at most $b_0$
and an initial slope in $[1/2,2]$.  Iterate the affine improvement using
\[
 U_i(x)=\frac{U(\lambda^i x)}{\lambda^i},
 \qquad b_i=b_0\lambda^{i\alpha_2}.
\]
Only constants are subtracted from the nonlinear unknown if a
renormalization of its value is needed; the affine functions are used
solely as comparison functions.  In particular the original
gradient-aligned equation remains its equation at every step.
The relative error threshold remains valid: geometric and forcing
errors acquire a factor $\lambda^i$, and
\[
 \lambda^i b_0\delta\leq b_i\delta
 \quad\text{because}\quad\alpha_2<1.
\]
The gradient remains bounded by its entry normalization, and the affine
slopes change by at most $Cb_i$.  Their total change is at most $1/8$,
so all slopes remain within $[1/4,4]$, as required by
Equation~\eqref{eq:elliptic:affine-improvement-hypotheses}.

Take $0<\alpha\leq\min(\alpha_d,\alpha_2)$.  Before affine entry the
constant approximations above have error $Cr^{1+\alpha}$.  At entry the
original radius is $r_e=r_*^{j+1}$ and the gradient normalization is
$a_e=k_*^j$.  At the subsequent radii $r=r_e\lambda^i$, the affine
errors in the original $u_0$ are bounded by
\[
 C a_e r_e\lambda^{i(1+\alpha_2)}
  =C r^{1+\alpha}\,
      \bigl(a_e r_e^{-\alpha}\bigr)
       \lambda^{i(\alpha_2-\alpha)}
  \leq C r^{1+\alpha},
\]
because
$a_e r_e^{-\alpha}=k_*^{-1}r_e^{\alpha_d-\alpha}\leq k_*^{-1}$.
Thus the constants remain uniform across every possible entry scale.
Intermediate radii are handled by restricting the approximation from
the next larger radius, whose ratio is one of the two fixed numbers
$r_*^{-1}$ and $\lambda^{-1}$.

We have proved that at every center and every sufficiently small radius
there is an affine approximation with uniform error $Cr^{1+\alpha}$.
The elementary affine Campanato criterion gives the gradient estimate:
comparison of the affine functions on concentric radii $r,r/2$ bounds
their slope difference by $Cr^\alpha$; summation identifies the limiting
slope with $Du$ at the center.  Comparing the approximations at two
centers in a common ball of radius comparable to their distance then
bounds the difference of these limiting slopes by that distance to the
power $\alpha$.  This proves the first estimate in
Equation~\eqref{eq:elliptic:rankone-conclusion}, also on the reflected
face, with constants uniform on the smaller original chart.

Finally let $L_x(y)=u(x)+Du(x)\cdot(y-x)$ in fixed coordinates.
The gradient estimate gives
$\|u-L_x\|_{L^\infty(B_{2r}(x))}\leq Cr^{1+\alpha}$.
Since
\[
 a:D^2(u-L_x)=G+a^{ij}\Gamma^\ell_{ij}u_\ell,
\]
Equation~\eqref{eq:elliptic:cordes-local} with $s=2$ yields
\[
 \|D^2u\|_{L^2(B_r(x))}
   \leq C\bigl(r^{n/2}+r^{n/2-1+\alpha}\bigr).
\]
The bounded connection term converting $D^2u$ to $\Hess_g u$ has
$L^2$ norm $O(r^{n/2})$.  Squaring the estimate and using $\alpha<1$
proves the Hessian estimate in
Equation~\eqref{eq:elliptic:rankone-conclusion}.  Restricting the reflected
estimate to the original half-ball proves the face version.
\end{proof}


\subsection{A bounded scalar equation with quadratic gradient growth}

The trace estimate now supplies the local mass control needed for
the second equation.  Its source has quadratic growth in $DZ$ and
also contains $|D^2f|^2$.  The latter need not yet have a pointwise
bound uniform over solutions: the estimate
$\int_{B_r}|D^2f|^2\le Cr^{n-2+2\alpha}$ from the preceding lemma
is the information available at this stage.

The following bounded-solution estimate uses exactly such a mass
bound.  Two exponential changes of variable absorb the quadratic
$DZ$ term, and the mass bound makes the remaining source potential
small on short scales.  Oscillation decay then gives a H\"older
modulus for $Z$ while its principal coefficients are still treated
as measurable.

\begin{lemma}[Natural-growth scalar estimate]
\label{lem:elliptic:natural-growth}
Let $n\ge3$ and let $Q$ be a Euclidean coordinate ball, or a smooth
flattened half-ball with constant Dirichlet value on its face.
Suppose a bounded classical function $Z$ satisfies weakly
\begin{equation}\label{eq:elliptic:natural-growth-model}
 \begin{gathered}
 -\partial_i(A^{ij}\partial_j Z)=\partial_i B^i+S,
 \qquad \lambda I\le A\le\Lambda I,\qquad A=A^{\mathsf T},\\
 |B|\le C_0,\qquad |S|\le C_0(1+|DZ|^2+F_0).
 \end{gathered}
\end{equation}
where $A$ is measurable and $F_0\ge0$ is bounded for each individual
solution.  Assume, for some $a\in(0,1)$ and $r_0>0$, and every ball
of radius $0<r\le r_0$ with center in the coordinate patch, that the measure
$\nu=(1+F_0)\,dx$ obeys
\begin{equation}\label{eq:elliptic:ng-source-morrey}
 \nu(B_r\cap Q)\le C_1r^{n-2+a}.
\end{equation}
Then $Z$ has a uniform $C^b$ bound on smaller patches, including the
constant-Dirichlet face, for some $b\in(0,a)$.  Its exponent and bound
depend only on $n,\lambda,\Lambda,C_0,C_1,a,r_0$, the bound for $|Z|$,
the fixed chart geometry and the distance from the remaining patch
boundary.  They use neither a modulus of continuity for $A$ nor a
uniform pointwise bound for $F_0$ or $D^2Z$.
\end{lemma}

\begin{proof}
If a Dirichlet face is present, first subtract its constant boundary
value.  This changes neither the equation nor its source bounds.
The reflection creates no measure on that face.  Flatten it as
$\{x_n=0\}$ and write $J=\operatorname{diag}(1,\ldots,1,-1)$.
For $x_n<0$ extend the four quantities by
\begin{equation}\label{eq:elliptic:ng-reflection}
 Z^e(x)=-Z(Jx),\quad A^e(x)=JA(Jx)J,\quad
 B^e(x)=-JB(Jx),\quad S^e(x)=-S(Jx).
\end{equation}
The zero trace makes $Z^e$ a $W^{1,2}$ function, and ellipticity is
preserved.  To verify the equation across the face, test the original
equation on the upper half with $\varphi(x)-\varphi(Jx)$, which has
zero trace.  Changing variables in the reflected half gives precisely
Equation~\eqref{eq:elliptic:natural-growth-model} for the extended quantities.  Equivalently,
the normal component of the full flux $ADZ+B$ agrees in the weak flux sense on the two sides.
Thus there is no boundary flux jump to estimate.  A different constant
Dirichlet value is handled by subtracting that constant first.  Reflect the nonnegative density $1+F_0$ evenly.  The ball-mass bound
changes by at most a fixed multiplicative factor.  The reflected
coefficients need not be smooth.

\medskip\noindent\emph{Exponential absorption and potentials.}
Use $\nu=(1+F_0)\,dx$, with the even extension just described where
needed.  We give the two-sign exponential argument for
Equation~\eqref{eq:elliptic:natural-growth-model}.  Put
$Q_\varepsilon=e^{\varepsilon LZ}$ for $\varepsilon\in\{-1,1\}$,
and let $\mathcal L=-\partial_i(A^{ij}\partial_j)$.
The weak chain rule gives
\begin{align}
 \mathcal L Q_\varepsilon
 &=\partial_i\bigl(\varepsilon LQ_\varepsilon B^i\bigr)
   -L^2Q_\varepsilon B\cdot DZ
   +\varepsilon LQ_\varepsilon S
   -L^2Q_\varepsilon A DZ\cdot DZ.\label{eq:elliptic:ng-exp-identity}
\end{align}
Indeed it follows directly by testing the equation for $Z$ with
$\varepsilon LQ_\varepsilon\varphi$.  Young's inequality bounds the
second term on the right by
$\tfrac12\lambda L^2Q_\varepsilon|DZ|^2+C L^2Q_\varepsilon$.
Taking $L\ge 2C/\lambda$ absorbs also the quadratic part of
$\varepsilon LQ_\varepsilon S$.  Since $Z$ is bounded and $L$ is
now fixed, we obtain, for both signs,
\begin{equation}\label{eq:elliptic:ng-exp-inequality}
 \mathcal L Q_\varepsilon
       \le \divg B_\varepsilon+C\nu,
 \qquad B_\varepsilon=\varepsilon LQ_\varepsilon B,
 \qquad \|B_\varepsilon\|_\infty\le C.
\end{equation}
This uses no derivative of $B$, $A$, or the density of $\nu$.

On a ball $B_r$ in the resulting full patch, let $v_\varepsilon$ solve
\begin{equation}\label{eq:elliptic:ng-potential}
 \mathcal L v_\varepsilon=\divg B_\varepsilon+C\nu,
 \qquad v_\varepsilon|_{\partial B_r}=0.
\end{equation}
The density of $\nu$ is bounded for each individual classical solution,
so this solve and the following tests can first be made in $W^{1,2}$.
Only the uniform mass bound in Equation~\eqref{eq:elliptic:ng-source-morrey} enters the
estimate.  The part with right side $\divg B_\varepsilon$ has absolute
value at most $Cr\|B_\varepsilon\|_\infty$.  This is the scaled
bounded-divergence-source estimate: testing by positive and negative
level truncations gives
$\int|D(v-k)_+|^2\le C\|B_\varepsilon\|_\infty^2
 |\{v>k\}|$, and Sobolev level iteration on the unit ball proves
the asserted bound after rescaling.

For the measure part, the scalar Dirichlet Green function of a
uniformly elliptic divergence operator in dimension $n\ge3$ satisfies
$0\le G_r(x,y)\le C|x-y|^{2-n}$.  For bounded measurable coefficients this follows from the
whole-space Green-function comparison
\cite[Section~7]{LittmanStampacchiaWeinberger1963}: extend $A$
uniformly elliptically off the ball, and compare the whole-space
and zero-Dirichlet Green functions by the maximum principle.  Dividing $B_r$ into annuli about any $x\in B_r$ gives
\[
 \int_{B_r}G_r(x,y)\,d\nu(y)
 \le C\sum_{j\ge0}(2^{-j}2r)^{2-n}
                   \nu\bigl(B_{2^{-j}2r}(x)\bigr)
 \le C\sum_{j\ge0}(2^{-j}2r)^a\le Cr^a.
\]
We work sufficiently far inside the coordinate patch that the mass
bound applies to the balls in this sum.  Consequently
\begin{equation}\label{eq:elliptic:ng-potential-bound}
 \|v_\varepsilon\|_{L^\infty(B_r)}\le E_r,
 \qquad E_r=C(r+r^a).
\end{equation}

\medskip\noindent\emph{Oscillation decay.}
Here is the oscillation reduction, including the role of the two
signs.  Let $m=\inf_{B_r}Z$, $M=\sup_{B_r}Z$, and $\omega=M-m$.
By Equations~\eqref{eq:elliptic:ng-exp-inequality} and~\eqref{eq:elliptic:ng-potential}, the functions
\[
 W_+=e^{LM}-Q_++v_++E_r,
 \qquad
 W_-=e^{-Lm}-Q_-+v_-+E_r
\]
are nonnegative weak supersolutions of $\mathcal L W\ge0$ on $B_r$.
At least one of the sets
\[
 \{Z\le(m+M)/2\}\cap B_{r/2},\qquad
 \{Z\ge(m+M)/2\}\cap B_{r/2}
\]
has at least half the measure of $B_{r/2}$.  On the first set
$e^{LM}-Q_+\ge c\omega$, and on the second
$e^{-Lm}-Q_-\ge c\omega$, because the exponential and its inverse
are uniformly Lipschitz on the fixed bounded range of $Z$.
For the corresponding sign the weak Harnack inequality, with its
fixed exponent $q>0$, yields
\[
 c\omega\le
 C\left(\frac1{|B_{r/2}|}\int_{B_{r/2}}W_\varepsilon^q\right)^{1/q}
 \le C\inf_{B_{r/4}}W_\varepsilon.
\]
Subtracting $v_\varepsilon+E_r$, whose absolute value is at most
$2E_r$, and using the same Lipschitz comparison with $Z$, improves
one of its two extremes.  Hence, for a fixed $\theta\in(0,1)$,
\begin{equation}\label{eq:elliptic:ng-oscillation-decay}
 \operatorname{osc}_{B_{r/4}}Z
 \le (1-\theta)\operatorname{osc}_{B_r}Z+C(r+r^a).
\end{equation}
Iteration gives a uniform $C^b$ bound for any fixed
\[
 0<b<\min\left\{a,\frac{-\log(1-\theta)}{\log4}\right\}.
\]
Applying the same conclusion to the reflected equation proves the
estimate up to the constant Dirichlet face.

\end{proof}

\subsection{Regularity of the coupled fixed points}

We apply the two local lemmas in a definite order.  The trace equation
first controls $Df$ and the local mass of $|D^2f|^2$.  The scalar
divergence lemma then controls the oscillation of $Z$.  These two
H\"older bounds permit scalar Schauder estimates, followed by an
absorption of the remaining quadratic gradient term.  Each step has
constants uniform in the outer truncation at fixed $N$.

\begin{proposition}\label{prop:elliptic:regularity}
Fix $N$.  Smooth solutions of either homotopy on $\mathcal M_{N,R}$
with $t\ge-\epsilon$ have bounds of every fixed order on coordinate
balls and on outer Dirichlet half-balls of a sufficiently small fixed
size.  Their size and the constants may depend arbitrarily on $N$,
but are independent of $R\ge R_{\min}(N)$ and of the homotopy
parameter.  In particular their $C^{1,\gamma}$ norms are bounded on
each fixed truncation, for any chosen $0<\gamma<1$.
\end{proposition}

\begin{proof}
\medskip\noindent\emph{1. The trace estimate and divergence form.}
Fix $N$ and a member of the homotopy.  Proposition~\ref{prop:elliptic:upper} gives
uniform bounds for $|f|$, $|df|$, and $|Z|$, and uniform ellipticity on
coordinate balls of a fixed small radius.  The radius and constants may
depend on $N$, but can be chosen independently of the outer truncation.
Lemma~\ref{lem:elliptic:gradient}, applied to the trace equation, gives some
$a\in(0,1)$ such that
\begin{equation}\label{eq:elliptic:ng-hessian-morrey}
 [df]_{C^a}\le C,
 \qquad
 \int_{B_r\cap\mathcal M_{N,R}}|D^2f|^2\,dx
       \le C r^{n-2+2a},\qquad 0<r\le r_0.
\end{equation}
Here and below the boundary version uses a smooth half-ball with constant
Dirichlet data; coordinate and covariant Hessians differ by a bounded
quantity.  Each occurrence of a ball in a mass estimate allows all
centers in the coordinate patch, not just one fixed center.

The implicit relation defining $t$ has derivative
$1+pv/k\ge1$ with respect to $t$.  On the bounded range now under
consideration it therefore defines a smooth function of $(x,Z,Df)$,
whose relevant partial derivatives are bounded.  In particular,
\begin{equation}\label{eq:elliptic:ng-dt}
 |Dt|\le C(1+|DZ|+|D^2f|).
\end{equation}
All the expressions in the equations are smooth before the block
decomposition, including at $df=0$.  Since $\mathcal T$ is quadratic
in $H^f$, $dt$, and $K$, Equation~\eqref{eq:elliptic:ng-dt} and the bounded
zeroth- and first-order quantities imply the following coordinate form
of the second equation:
\begin{equation}\label{eq:elliptic:ng-natural-equation}
 -\partial_i(A^{ij}\partial_j Z)=\partial_i B^i+S,
 \qquad
 \lambda I\le A\le\Lambda I,
 \qquad |B|\le C,
 \qquad |S|\le C\bigl(1+|DZ|^2+|D^2f|^2\bigr).
\end{equation}
The matrix $A$ is symmetric.  With volume densities included,
$A$ represents $kuA_\chi$, $B$ represents $uA_\chi K(w,\cdot)$,
and $S=-\sqrt{\det g}\,u\Xi$.  The same bounds hold uniformly along
both stages of the homotopy; scaling $\Xi$ towards zero improves its
absolute-value bound.  No continuity modulus for $A$ is being used yet.

Apply Lemma~\ref{lem:elliptic:natural-growth} with
$F_0=|D^2f|^2$.  Its coefficient and source hypotheses are precisely
Equation~\eqref{eq:elliptic:ng-natural-equation}.  On patches of size
at most one, Equation~\eqref{eq:elliptic:ng-hessian-morrey} gives
$\nu(B_r)\le Cr^{n-2+a}$ for
$\nu=(1+|D^2f|^2)\,dx$.  The lemma therefore gives a uniform
$C^b$ bound for $Z$ for some $0<b<a$, up to the constant outer
Dirichlet face as well.  Both $Z$ and $Df$ now have the continuity
needed to improve the trace equation by Schauder theory.

\medskip\noindent\emph{2. Absorbing the quadratic gradient term.}
The coefficients of the trace equation are smooth functions of
$(x,Z,Df)$.  The preceding estimate and Equation~\eqref{eq:elliptic:ng-hessian-morrey}
therefore give a uniform $C^c$ modulus for those coefficients and
the trace right side, with $c=\min(a,b)>0$.  Interior and Dirichlet
Schauder estimates give
\begin{equation}\label{eq:elliptic:ng-f-schauder}
 \|f\|_{C^{2,c}}\le C
\end{equation}
on smaller patches of the same uniform type.  This use of Schauder
requires only the now established H\"older modulus of $(Z,Df)$.

Expand the divergence equation on each original ball or half-ball.
Since $A=A(x,Z,Df)$ and $B=B(x,Z,Df)$, their derivatives consist of
bounded terms, terms linear in $DZ$, and terms linear in $D^2f$.
Equation~\eqref{eq:elliptic:ng-f-schauder} and the formula for $Dt$ show that
the expanded equation has the form
\begin{equation}\label{eq:elliptic:ng-expanded}
 A^{ij}D_{ij}Z=H,
 \qquad |H|\le C(1+|DZ|^2),
 \qquad \|A\|_{C^c}\le C.
\end{equation}
The bound on the H\"older norm of $A$ does not presuppose a bound for
$DZ$.  Coordinate lower-order terms are included in $H$.

Fix any finite $s>n$.  Freezing the now H\"older principal matrix
on sufficiently small balls gives the interior and flat-Dirichlet
$W^{2,s}$ estimates with uniformly bounded coefficient in front of
$\|H\|_s$.  On nested patches $Q_r\subset Q_{2r}$ these give
\begin{equation}\label{eq:elliptic:ng-local-w2s}
 \|D^2Z\|_{s,Q_r}
 \le C_s\|DZ\|_{2s,Q_{2r}}^2+C_{s,r}.
\end{equation}
The last constant includes the bounded height, the constant term
in Equation~\eqref{eq:elliptic:ng-expanded}, and the scaled cutoff costs.
For a ball or a fixed-shape half-ball the scaled interpolation
inequality is
\begin{equation}\label{eq:elliptic:ng-interpolation}
 \|DZ\|_{2s,Q}^2
 \le C_s\operatorname{osc}_Q Z\,\|D^2Z\|_{s,Q}
   +C_s r^{n/s-2}(\operatorname{osc}_Q Z)^2,
\end{equation}
where $Q$ has size $r$.  Indeed apply the usual second-order
interpolation inequality on the unit ball or half-ball after
subtracting a constant and rescaling.  On a half-ball, the constant
Dirichlet value belongs to the closed range of $Z$, so it can be
used for this subtraction.  The constants depend only on the
uniform chart shapes, $s$, and the dimension.

For clarity, the absorption is local and does not lose the number
of patches in a large truncation.  For a fixed sufficiently small
$r$, let
\[
 \mathcal M_r=\sup_x\|D^2Z\|_{s,Q_r(x)}.
\]
It is finite for each individual classical solution on a truncation.
The uniform local geometry allows each $Q_{2r}(x)$ to be covered
by at most $J_0$ patches of size $r$, with $J_0$ independent of the
outer radius.  Thus Equations~\eqref{eq:elliptic:ng-local-w2s} and
\eqref{eq:elliptic:ng-interpolation} imply
\[
 \mathcal M_r
 \le C_s J_0^{1/s}
       \sup_x\operatorname{osc}_{Q_{2r}(x)}Z\,\mathcal M_r
       +C_{s,r}.
\]
Choose $r$ small enough, using the uniform $C^b$ bound already
proved, that the displayed coefficient of $\mathcal M_r$ is at
most $1/2$.  This gives a uniform local $W^{2,s}$ bound up to the
outer face, independently of the total number of patches.  Sobolev
embedding now bounds $Z$ in $C^{1,d}$ for $d<1-n/s$.

\medskip\noindent\emph{3. Bootstrap.}
The trace equation consequently bounds $f$ in $C^{3,d'}$ for some
$d'>0$.  Its coefficients now have one H\"older derivative; the
expanded second equation has H\"older right side, giving
$Z\in C^{2,d'}$.  Repetition yields bounds of every fixed order,
since the background data are smooth.  More explicitly, once
$Z\in C^{j,d'}$ and $f\in C^{j+1,d'}$ for $j\ge1$, the trace
equation first gives $f\in C^{j+2,d'}$, and the divergence equation
then gives $Z\in C^{j+1,d'}$.  Shrinking the exponent once at the
first step suffices.  The argument applies initially to the
$C^{2,\gamma}$ fixed points constructed below: their individual
Hessian norms are already finite, and smoothness is the conclusion
of the bootstrap.  All constants in this paragraph may depend
arbitrarily on fixed $N$ and the chosen derivative order.

\end{proof}

\subsection{The trace equation for an arbitrary trial}

The degree argument requires a map defined beyond the region $t>-\epsilon$.
The next lemma supplies it.  Its estimates are separate from
Proposition~\ref{prop:elliptic:upper}: here a trial $Z$ is prescribed,
and no lower bound for its associated $t$ is assumed.

\begin{lemma}\label{lem:elliptic:trial-trace}
Let $\mathcal D$ be a compact smooth $n$-dimensional Riemannian
manifold with nonempty smooth boundary, with $n\ge3$ and $k=n-1$,
and let $K$ be a smooth symmetric tensor on $\mathcal D$.
Fix $N>k$, $0<\gamma<1$, and a positive constant $\eta$.
Use $p(t)=N\vartheta(Nt)$ and
$l(t)=\exp(\int_0^t p(s)\,ds)$, with the smooth cutoff
$\vartheta$ of Equation~\eqref{eq:scalar:parameters}, and define
$D$, $t$, $A_\chi$ by Equations~\eqref{eq:scalar:implicit} and
\eqref{eq:scalar:variables}.  Put $K_a=aK$.
For every $z\in C^{1,\gamma}(\mathcal D)$ and $a\in[0,1]$,
the Dirichlet problem
\[
 \tr_{A_\chi}\left(K_a+\frac l{\sqrt D}\Hess f\right)=\eta f,
 \qquad f|_{\partial\mathcal D}=0,
 \qquad t+\frac1{2k}\log(1+l(t)^2|df|^2)=z
\]
has a unique solution $f=f_a[z]\in C^{3,\gamma}$.
The solution map is continuous from
$[0,1]\times C^{1,\gamma}(\mathcal D)$ to
$C^{3,\gamma}(\mathcal D)$ and maps bounded sets to bounded sets.
All constants may depend on $N,\eta,\mathcal D,g,K$ and the
prescribed bound for $z$.  There is no restriction on the minimum
of the implicitly determined $t$.
\end{lemma}

\begin{proof}
\medskip\noindent\emph{1. Coefficient growth without a floor.}
Let $b=\sqrt D/l$.  On a bounded trial range the coefficients satisfy,
for every $\sigma=|df|$,
\begin{equation}\label{eq:elliptic:trial-growth}
 b\le C(1+\sigma),\qquad
 \frac{c}{1+\sigma}\le\chi\le1.
\end{equation}
We verify the bounds also at large slope, where $t$ has no fixed floor.
The implicit relation implies $t<0$ when $\sigma$ is sufficiently large.
There $p=N$, $l=e^{Nt}$, and
\[
 e^{2kz}=e^{2kt}+\sigma^2e^{2(k+N)t}.
\]
Uniformly on a bounded $z$ range, as $\sigma\to\infty$,
\[
 t=\frac{kz-\log\sigma}{k+N}+o(1),\qquad
 d\asymp\sigma^{-2k/(k+N)},\qquad
 \chi\asymp\frac{k}{k+N}\sigma^{-2k/(k+N)}.
\]
The exponent is less than one because $N>k$.
Moreover $b=\sigma/\sqrt v\asymp\sigma$ there.
The remaining bounded slope range is compact in $(z,\sigma)$, and
the coefficients are smooth and positive on it.  This proves
Equation~\eqref{eq:elliptic:trial-growth}.

\medskip\noindent\emph{2. A global gradient bound.}
Continue from $\Delta_g f=f$ by solving, for $0\le q\le1$,
\begin{equation}\label{eq:elliptic:trace-continuation}
 A_q:\Hess f-c_q f=-q b\tr_{A_\chi}K_a,
 \qquad A_q=(1-q)I+qA_\chi,
 \qquad c_q=1-q+q\eta b.
\end{equation}
The scalar $c_q$ is positive.  Extrema give a common height bound
$\|f\|_\infty\le H$, since
\[
 \frac{q b|\tr_{A_\chi}K_a|}{c_q}\le\frac C\eta.
\]
On this height range the right side of the equation written as
$A_q:\Hess f=G(x,f,df)$ satisfies, for a covector $\xi$,
\begin{equation}\label{eq:elliptic:trial-forcing}
 |G(x,f,\xi)|\le C(1+|\xi|),
\end{equation}
uniformly in $q,a$.  The eigenvalues of $A_q$ are one in the
transverse directions and $\chi_q=1-q+q\chi\ge c/(1+|\xi|)$ along $\xi$.

Here is a gradient estimate using precisely this lower axial bound.
Choose $K_0$ sufficiently large, then a positive distance $d_0$ smaller
than the boundary tubular radius, the injectivity radius of a smooth
extension of the compact domain, and $(2K_0)^{-1}$.
For $M>0$ set
\begin{equation}\label{eq:elliptic:log-modulus}
 \psi(d)=K_0^{-1}\log(1+K_0Md),\qquad
 \psi''=-K_0(\psi')^2.
\end{equation}
By increasing $M$ we ensure both $\psi'\ge1$ on $[0,d_0]$ and
$\psi(d_0)>2H$.  Indeed $\psi'(d_0)$ tends to
$(K_0d_0)^{-1}>2$, while $\psi(d_0)$ tends to infinity.

First compare with $\pm\psi(d(x,\partial\mathcal D))$ in the boundary collar.
At a contact, the gradient is normal and has length $\sigma=\psi'$.
The tangential Hessian contraction costs at most $C\sigma$, whereas
\[
 \chi_q\psi''\le-\frac{cK_0\sigma^2}{1+\sigma}
                         \le-\frac{cK_0}2\sigma.
\]
Choose $K_0$ large enough to dominate the collar curvature and
Equation~\eqref{eq:elliptic:trial-forcing}.  The positive barrier is a
strict supersolution; the negative one is a strict subsolution.  On
the inner edge of this collar they dominate $\pm f$ by the height
bound, and both vanish at the boundary.  Therefore
\begin{equation}\label{eq:elliptic:boundary-modulus}
 |f(x)|\le\psi(d(x,\partial\mathcal D))
 \quad\hbox{when }d(x,\partial\mathcal D)\le d_0.
\end{equation}

Consider now $f(x)-f(y)-\psi(d(x,y))$ for pairs with
$0\le d(x,y)\le d_0$, using the distance of the smooth ambient
extension.  It is nonpositive on the diagonal, for $d=d_0$, and when
either endpoint lies on the boundary, by
Equation~\eqref{eq:elliptic:boundary-modulus} and the monotonicity of
$\psi$.  Suppose it had a positive interior maximum.  The endpoint
gradients are parallel along their unique short joining geodesic and
have the same length $\sigma=\psi'\ge1$.  Write $e$ for their common
parallel direction.  Varying one endpoint in the axial direction
gives
\[
 \Hess f_x(e,e)\le\psi'',\qquad
 \Hess f_y(e,e)\ge-\psi''.
\]
Varying both endpoints in parallel transverse directions and summing
the second-variation inequalities gives
\[
 \tr_{e^\perp}\Hess f_x-
 \tr_{e^\perp}\Hess f_y\le C d(x,y)\sigma.
\]
For completeness, the simultaneous endpoint variation is generated
by a parallel vector field along the geodesic; its index form is
bounded above by $C d(x,y)$ since the ambient curvature is bounded.
This is the displayed estimate after multiplication by $\psi'$.
The transverse eigenvalues are exactly one at both endpoints, so
\begin{align*}
 A_{q,x}:\Hess f_x-A_{q,y}:\Hess f_y
 &\le Cd(x,y)\sigma+(\chi_{q,x}+\chi_{q,y})\psi''\\
 &\le Cd_0\sigma-cK_0\sigma.
\end{align*}
The right sides of Equation~\eqref{eq:elliptic:trace-continuation}
differ by a quantity at least $-2C(1+\sigma)$.
A sufficiently large $K_0$ contradicts this.  Notice that no estimate
for $\chi_{q,x}-\chi_{q,y}$ was needed: both signed axial Hessian
bounds have the favorable sign.  We conclude
$|f(x)-f(y)|\le\psi(d(x,y))$ and hence $|df|\le\psi'(0)=M$.

\medskip\noindent\emph{3. Continuation and regularity.}
Equation~\eqref{eq:elliptic:trial-growth} now gives uniform ellipticity
throughout the continuation.  Its normalized forcing is bounded.
Lemma~\ref{lem:elliptic:gradient} gives a common positive H\"older
exponent for $df$, including the boundary.  Since the coefficients
are smooth functions of $(x,z,df)$, the Dirichlet Schauder estimate
first gives $f\in C^{2,\gamma_0}$ for a possibly smaller exponent.
The resulting Lipschitz bound for $df$ then makes those coefficients
$C^{0,\gamma}$, so the estimate improves to $C^{2,\gamma}$.
Differentiating the equation once gives $C^{3,\gamma}$ estimates:
the trial is $C^{1,\gamma}$ and all differentiated coefficient and
source terms are controlled.  These scalar estimates are uniform
along the continuation; they are the interior and smooth Dirichlet
estimates of~\cite{GilbargTrudinger} on a finite coordinate covering.

The linearization in $f$ is a uniformly elliptic scalar operator with
bounded H\"older first-order coefficients and zero-order coefficient
$-c_q<0$.  Its homogeneous Dirichlet kernel is zero by the maximum
principle.  The scalar linear Dirichlet continuity theorem, starting
from $\Delta_g-1$ and using the Schauder estimate, makes it invertible
from $C^{2,\gamma}_0$ to $C^{0,\gamma}$.
The implicit function theorem gives openness of the nonlinear
continuation set; the bounds and compactness give closedness.
The initial solution is $f=0$, hence existence follows at $q=1$.

\medskip\noindent\emph{4. Uniqueness and continuity.}
For uniqueness compare two solutions at a positive maximum of their
difference.  Their gradients agree there, so their coefficient
matrices and their positive factors $b$ agree; Hessian ordering and
the strictly negative zero-order term give a contradiction.
For continuity in the stated norm, regard the normalized trace
operator as a map
\[
 [0,1]\times C^{1,\gamma}\times C^{3,\gamma}_0
                    \longrightarrow C^{1,\gamma}.
\]
It is continuously differentiable in the trial and solution variables,
since its coefficient functions are smooth on every bounded jet range.
At a solution its linearized principal, first-order, and zero-order
coefficients are $C^{1,\gamma}$.  The same maximum principle and
the one-derivative-higher Dirichlet Schauder estimate therefore make
that linearization an isomorphism from $C^{3,\gamma}_0$ onto
$C^{1,\gamma}$.  The implicit function theorem supplies local
continuous dependence into $C^{3,\gamma}$; uniqueness joins these
local maps.  The a priori estimates already proved give the
bounded-set assertion.
\end{proof}

\subsection{A compact map and the two homotopies}

We now return to $\mathcal D=\mathcal M_{N,R}$ and
$\eta=\ell^{3/2}$.  Fix $N,R$ with $R\ge R_{\min}(N)$, and let
\[
 X=\{z\in C^{1,\gamma}(\mathcal M_{N,R}):z|_{S_R}=0\},
 \qquad 0<\gamma<1.
\]
For a trial $z\in X$, solve $f=f_a[z]$ by
Lemma~\ref{lem:elliptic:trial-trace}.  Evaluate $t,D,u,A_\chi,w$ and
$\mathcal T$ using $(f,z)$, including the derivatives of $t$ obtained
by differentiating its implicit definition.  Define $T_a(z)=Z_{\rm out}$
to be the solution of the \emph{linear} Dirichlet problem
\begin{equation}\label{eq:elliptic:compact-map}
 \divg(ku A_\chi\nabla Z_{\rm out})
   =u\Xi_a-\divg\bigl(uA_\chi K_a(w,\cdot)\bigr),
 \qquad Z_{\rm out}|_{S_R}=0.
\end{equation}
Only the displayed gradient of $Z_{\rm out}$ is unfrozen.
The source $\Xi_a$ is evaluated at the input $z$, not at the output.

On bounded subsets of $X$, the trace lemma gives bounded
$C^{3,\gamma}$ norms for $f$.  Consequently the principal coefficient
$kuA_\chi$ is $C^{1,\gamma}$ and uniformly positive definite, the
drift flux is $C^{1,\gamma}$, and $u\Xi_a$ is $C^{0,\gamma}$.
The linear scalar Dirichlet theorem and Schauder estimate therefore
give a unique $C^{2,\gamma}$ output.  The map $T_a:X\to X$ is continuous
and compact, uniformly on bounded sets and $a\in[0,1]$.
The same construction, with $a=0$ and source $s u\Xi_0$, defines
the second compact homotopy $T_{0,s}$.  Its trace solution is zero for
every trial, since $K=0$.

Every fixed point bootstraps to smoothness.  Indeed the first
Schauder step just described gives $z\in C^{2,\gamma}$; the trace
equation then gives $f\in C^{4,\gamma}$, and repeated application to
the second equation and then the first gives arbitrary derivatives.
Thus the a priori estimates proved above apply to all fixed points.

For either homotopy parameter $\lambda$, set
\begin{equation}\label{eq:elliptic:degree-domain}
 \mathcal U_\lambda
 =\{z\in X:\min_{\mathcal M_{N,R}}t[z,df_\lambda[z]]>-\epsilon,
                         \ \|z\|_{C^{1,\gamma}}<M\}.
\end{equation}
The associated set of parameter--trial pairs is
\[
 \mathcal U=\{(\lambda,z)\in[0,1]\times X:z\in\mathcal U_\lambda\}.
\]
It is relatively open in $[0,1]\times X$ because the trace solution
map and the implicit function defining $t$ are continuous.  Choose $M$ larger than the common $C^{1,\gamma}$ bound
for all above-floor fixed points supplied by
Proposition~\ref{prop:elliptic:regularity}.  There is no fixed point
on the norm boundary, and the floor lemmas exclude one on the other
boundary.

Here is the precise reason degree may be used on these moving sets.
The fixed points in the closure of $\mathcal U$ form a compact subset of
$[0,1]\times X$: the parameter interval is compact, the trial norms
are bounded, and compactness of the homotopy gives a convergent
subsequence of every sequence of such fixed points.  Its limit cannot lie on either excluded boundary.  Near any one parameter, take a fixed open
neighborhood of that compact fiber whose closure remains inside
$\mathcal U_\lambda$ for all sufficiently nearby parameters; shrink
the parameter interval to exclude other fixed points outside this
neighborhood, using compactness.  Excision and ordinary fixed-domain
homotopy invariance identify the degrees at those nearby parameters.
A finite chain of such intervals identifies the degree at the two
ends.  These are the usual excision and homotopy properties of
Leray--Schauder degree~\cite[Chapter~2]{Deimling1985}.

At the endpoint $a=0,s=0$, Equation~\eqref{eq:elliptic:compact-map}
has zero source and zero boundary value, so $T_{0,0}(z)=0$ for every
trial.  Its sole fixed point is zero, with $t=0>-\epsilon$, and its
degree is one.  First follow the second homotopy to $s=1$, then the
first to $a=1$.  The degree remains one, yielding a smooth solution
of the desired filled system on every sufficiently large truncation.

\subsection{Exhaustion, end decay, and the energy flux}

The degree argument gives solutions on every sufficiently large
truncation.  Uniform local estimates let the outer boundary escape
while $N$ remains fixed.  The positive zeroth-order term in the trace
equation then gives exponential decay of the graph; a scalar change
of variable converts the end equation for $Z$ to a Poisson estimate.
These are the precise asymptotics needed to evaluate the energy.

\begin{proof}[Completion of the proof of Theorem~\ref{thm:elliptic:filled}]
Keep $N$ fixed and let $R\to\infty$.  The estimates of
Proposition~\ref{prop:elliptic:regularity} on compact sets are
independent of $R$.  A diagonal subsequence therefore converges
smoothly on every compact subset of $\mathcal M_N$ to a solution
$(f,Z)$.  The polynomial bounds and height comparisons survive.
They give $t\ge-\epsilon$; equality would be an interior global
minimum and contradict Lemma~\ref{lem:scalar:floor}.

We prove the remaining end estimates on the truncation solutions,
uniformly in $R$ at this fixed $N$; the estimates then pass to the
smooth limit just obtained.

On the end $K=0$ and the trace equation reads
\begin{equation}\label{eq:elliptic:end-trace}
 A_\chi:\Hess f-\frac{\eta\sqrt D}{l}f=0,
 \qquad \frac{\eta\sqrt D}{l}\ge\frac\eta e>0.
\end{equation}
For small $a_N>0$ and a sufficiently large fixed $r_2(N)$,
$v_N(r)=e^{-a_N(r-r_2)}$ obeys
\[
 A_\chi:\Hess v_N-\frac{\eta\sqrt D}{l}v_N<0
 \quad(r\ge r_2).
\]
In fact $0<A_\chi\le I$ and the end metric bounds give
$A_\chi:\Hess v_N\le C(a_N^2+a_N/r)v_N$; choose
$Ca_N^2<\eta/(4e)$ and then $r_2$ large.  Multiples of $\pm v_N$
dominate $f$ at $S_{r_2}$ and at the outer zero sphere.  The maximum
principle for the operator in Equation~\eqref{eq:elliptic:end-trace},
with its coefficients evaluated at the solution, proves exponential
decay uniformly on all truncations.  The already established scalar
local estimates, on balls and boundary half-balls of a fixed size,
then give the same decay for every derivative of $f$.
In particular $t-Z$ and its derivatives are exponentially small.

The second equation can consequently be written, with an error
exponentially small through every fixed number of derivatives, as
\begin{equation}\label{eq:elliptic:end-equation}
\begin{gathered}
 k\Delta_g Z+k B_N(Z)|dZ|_g^2
 =\rho-\vartheta(N(Z+\epsilon))C_N\rho_0+O_j(e^{-a_Nr}),\\
 B_N(z)=(k-1)+p(z)-\delta_0s_{p(z)}.
\end{gathered}
\end{equation}
To verify this formula, set $f=0$ in the coefficients.  Then
$D=1$, $t=Z$, $w=0$, $A_\chi=I$,
$\mathcal T=ks_{p(Z)}|dZ|^2$, and
$V=ku(Z)\nabla Z$ with $(\log u)'=k-1+p$.
Every difference from these expressions has an exponentially small
factor involving a derivative of $f$; its other factors are bounded
at fixed $N$ by local regularity.  This also proves the differentiated
error assertion.  The outer boundary causes no loss, since the same
fixed-size local estimates hold there uniformly in $R$.

Define a smooth increasing function $Q_N$ on the bounded range of $Z$
by
\[
 Q_N(0)=0,\qquad Q_N'(z)
 =\exp\left(\int_0^z B_N(\tau)\,d\tau\right).
\]
Its derivative has positive upper and lower bounds at fixed $N$.
Equation~\eqref{eq:elliptic:end-equation} and the chain rule give
\begin{equation}\label{eq:elliptic:end-poisson}
 |\Delta_g Q_N(Z)|\le C_N' r^{-n-\beta},
 \qquad Q_N(Z)|_{S_R}=0.
\end{equation}
Here $C_N'$ denotes a constant at fixed $N$, not necessarily the
penalty polynomial.  The function
\[
 b_*(r)=r^{2-n}(1-r^{-\beta/2})
\]
is positive for $r>1$ and satisfies
\[
 \Delta_g b_*
 =-\frac\beta2\left(n-2+\frac\beta2\right)r^{-n-\beta/2}
          +O(r^{2-2n})<0
\]
sufficiently far out.  The correction dominates the metric error
because $\beta/2<n-2$.
Taking a sufficiently large multiple of $b_*$ dominates both the
source in Equation~\eqref{eq:elliptic:end-poisson} and the inner
boundary values.  Comparison with its two signs, also valid on the
outer zero sphere, gives
$|Q_N(Z)|\le C_N' r^{2-n}$ on every truncation.  Monotonicity of $Q_N$
gives $Z=O(r^{2-n})$, and hence the same bound for $t$.

Rescale successive end annuli to unit size.  The metric coefficients
have uniformly controlled derivatives and converge to the Euclidean
ones.  The right side of the equation for $Q_N(Z)$ is a smooth
function of $Z$ times $r^{-n-\beta}$ plus exponentially small terms.
Interior Poisson $W^{2,p}$ estimates, then Schauder estimates and
differentiation, give $Z,t=O_j(r^{2-n})$ for each $j$; the constants
may depend on $N,j$.  Finally
$|dZ|^2=O(r^{2-2n})=O(r^{-n-\beta})$, since $\beta<n-2$.
Equation~\eqref{eq:elliptic:end-equation} now proves the stated
$\Delta_g t$ bound and finishes the theorem.
\end{proof}

For the resulting solution, Proposition~\ref{prop:scalar:flux} now
applies with all of its end hypotheses verified.  In particular,
writing $\omega=\omega_{n-1}$ and
$\mathfrak F_N=\lim_{r\to\infty}\int_{S_r}g(V,\nu_g)\,dA_g$,
\begin{equation}\label{eq:elliptic:constructed-energy}
 \widehat E-E_0=-\frac{\mathfrak F_N}{k\omega},\qquad
 \mathfrak F_N=\int_{\mathcal M_N}u\Xi\,dV_g
       \ge-\mathcal P_N(\ell+\ell^{1/2}).
\end{equation}
Here $E_0$ is the energy of the fixed prepared data.  The metric also
satisfies $\widehat g\ge e^{-2\epsilon}g$ pointwise.  The flux estimate
uses the polynomial penalty-band calculation of that proposition;
the arbitrarily $N$-dependent constants in the end regularity estimates
are used only to justify the flux and its normalization.

\begin{remark}\label{rem:elliptic:limit-order}
The construction exhausts the radii only after $N$ has been fixed.
Equations~\eqref{eq:elliptic:polynomial-bounds} and
\eqref{eq:elliptic:height-bound} are the estimates used with $N\to\infty$
in the next section.  No bound for a high derivative that depends
arbitrarily on $N$ is substituted for a polynomial bound.  In
particular, this proof takes no elliptic limit as $N\to\infty$ and
imposes no equation on a limiting singular hypersurface.
\end{remark}
